\subsection{两个线性映射的张量积}

设A是一个环，E、E'是两个右A-模，F、F'是两个左A-模，且 \(u: E \rightarrow E'\) 和 \(v: F \rightarrow F'\) 是两个A-线性映射。容易验证，映射

\[
(x, y) \mapsto u (x) \otimes v (y)
\]

（从 \(\mathbf{E} \times \mathbf{F}\) 到 \(\mathbf{E}' \otimes_{\mathbf{A}} \mathbf{F}'\) ）是 \(\mathbf{Z}\) -双线性的，且满足第1号的条件(1)。因此，根据第1号命题1，存在且仅存在一个 \(\mathbf{Z}\) -线性映射 \(w: \mathbf{E} \otimes_{\mathbf{A}} \mathbf{F} \to \mathbf{E}' \otimes_{\mathbf{A}} \mathbf{F}'\) ，使得

\[
w (x \otimes y) = u (x) \otimes v (y)\tag{3}
\]

对于 \(x \in \mathbf{E}\) ， \(y \in \mathbf{F}\) 。此映射记作 \(u \otimes v\) （当不致引起混淆时），并称为线性映射 \(u\) 与 \(v\) 的张量积。

由(3)立即得出 \((u, v) \mapsto u \otimes v\) 是一个Z-双线性映射，称为典范映射

\[
\operatorname{Hom} _ {\mathbf {A}} (\mathrm{E}, \mathrm{E} ^ {\prime}) \times \operatorname{Hom} _ {\mathbf {A}} (\mathrm{F}, \mathrm{F} ^ {\prime}) \rightarrow \operatorname{Hom} _ {\mathbf {Z}} (\mathrm{E} \otimes_ {\mathbf {A}} \mathrm{F}, \mathrm{E} ^ {\prime} \otimes_ {\mathbf {A}} \mathrm{F} ^ {\prime}).
\]

根据第1号命题1，与之对应的是一个 \(\mathbf{Z}\) -线性映射，称为典范映射

\[
\operatorname{Hom} _ {\mathbf {A}} (\mathrm{E}, \mathrm{E} ^ {\prime}) \otimes_ {\mathbf {Z}} \operatorname{Hom} _ {\mathbf {A}} (\mathrm{F}, \mathrm{F} ^ {\prime}) \rightarrow \operatorname{Hom} _ {\mathbf {Z}} (\mathrm{E} \otimes_ {\mathbf {A}} \mathrm{F}, \mathrm{E} ^ {\prime} \otimes_ {\mathbf {A}} \mathrm{F} ^ {\prime})\tag{4}
\]

它把张量积的每个元素 \(u \otimes v\) 对应到线性映射 \(u \otimes v: \mathbf{E} \otimes_{\mathbf{A}} \mathbf{F} \to \mathbf{E}' \otimes_{\mathbf{A}} \mathbf{F}'\) 。注意，典范映射(4)不一定是单射的，也不一定是满射的。因此记号 \(u \otimes v\) 可能引起混淆，须由上下文指明它表示的是张量积的元素还是线性映射。

此外，设 \(\mathbf{E}''\) 是一个右A-模， \(\mathbf{F}''\) 是一个左A-模，且 \(u':\mathbf{E}'\to\mathbf{E}''\) , \(v':\mathbf{F}'\to\mathbf{F}''\) 是A-线性映射；由(3)得出

\[
\left(u ^ {\prime} \circ u\right) \otimes \left(v ^ {\prime} \circ v\right) = \left(u ^ {\prime} \otimes v ^ {\prime}\right) \circ (u \otimes v).\tag{5}
\]
% semantic-fragments: legacy-input-seam
\relax{}
